\section{合卷：西周留下的，不只是“分封制”三个字}

若把西周写成“分封，所以后来诸侯强大”，历史会变得太扁。更准确的说法是：周用一套兼具宗族关系、礼仪名分、资源赏赐和军事义务的网络，解决了征服之后如何低成本统治的问题。它在很长时间里有效；正因为有效，地方也得以在网络内部积累能力。

西周的得失不在于它是否“先进”，而在于它面对的条件：中心行政能力有限，空间广阔，旧商与新周的关系需要被重新安置。它的成功是把胜利变成秩序；它的失败是当继承、财政、边疆和内部信任同时紧张时，没有足够直接的制度工具把网络重新拢回中心。

这种遗产不会在公元前771年消失。东周诸侯仍要借“周王室”的名分说话，春秋霸主仍要在会盟中使用旧秩序的语言；只是实际资源已经越来越掌握在诸侯手里。这正是下一章的起点。
